\scanpage{8}
% archive foliation: 7
\deleted{revient à la donnée d'un scindage
\diagram{\begin{tikzcd}[ampersand replacement=\&]
\mathcal G'\ar[r]\&\Gamma\ar[l,bend left=35]
\end{tikzcd}}}

\textcircled{2.3.} Soit
\[
(17)\qquad\mathbb G'=\operatorname{Aut}_{\bar k}(\mathcal X'\longrightarrow\mathcal X)
\]
\deleted{\unsure{donc} $\mathbb G'$ \unsure{se dé}} donc suite exacte
\[
(18)\quad\diagram{\begin{tikzcd}[ampersand replacement=\&]
1\ar[r]\&G\ar[r]\&\mathbb G'\ar[r]\&\mathbb G.
\end{tikzcd}},
\]
qui se déduit (\unsure{mod. iso. can.}) de (16) comme
\[
\begin{aligned}
(19)\qquad\mathbb G'&\simeq\text{image inverse de }\mathbb G\subset\mathcal G\\
&\qquad\text{par }\mathcal G'\longrightarrow\mathcal G\\
&\simeq\operatorname{Ker}(\mathcal G'\longrightarrow\Gamma)
\end{aligned}
\]

Le cas favorable $N$ inv. dans $\mathcal E$ i.e. $\mathcal G'\to\mathcal G$ surjectif revient à dire que a) $\mathbb G'\to\mathbb G$ surjectif et b) $\mathcal G'\to\Gamma$ \deleted{surjectif} \unsure{même image} que $\mathcal G\to\Gamma$.

\deleted{\uline{NB}}

\uline{NB} Si $\mathcal X$ de prés. finie sur $\unsure{\xi|\bar k}$, alors l'image $\Gamma_1$ de $\mathcal G$ dans $\Gamma$ est un sous-groupe ouvert, et si de plus $\mathcal X'$ est fini sur $\mathcal X$ i.e. de prés. finie sur $k$, alors l'image $\Gamma'_1\inserted{(\subset\Gamma_1)}$ de $\mathcal G'$ dans $\Gamma$ l'est également.
